\documentclass[12pt]{beamer}
\usetheme{moloch}

\usepackage{amsmath,amscd}
\usepackage{fontspec}
\usepackage{unicode-math}
\usepackage[dvipsnames]{xcolor}

\setbeamercolor{background canvas}{bg=black}
\setbeamercolor{normal text}{fg=white}

\setmonofont{DejaVu Sans Mono}

% Set the title bar to black background with white text
\setbeamercolor{frametitle}{bg=black, fg=white}

% Ensure other structural elements (bullets, etc.) are visible
\setbeamercolor{structure}{fg=white}

\setmainfont{Libertinus Serif}
\setsansfont{Libertinus Sans}
\setmonofont{DejaVu Sans Mono}
\setmathfont{STIXTwoMath-Regular.otf}
%\setmathfont{STIX Two Math}


% Define colors for syntax highlighting
\definecolor{codegreen}{rgb}{0,0.6,0}
\definecolor{codegray}{rgb}{0.5,0.5,0.5}
\definecolor{codecoffee}{rgb}{0.75,0.6,0.3}
\definecolor{question}{rgb}{1.0,0.75,0.1}
\definecolor{lightblue}{rgb}{0.55,0.75,0.9}
\definecolor{yellow}{rgb}{1.0, 1.0, 0.0}
\definecolor{codecomment}{rgb}{0.97, 0.74, 0.4}
\definecolor{codepy}{rgb}{0.9, 0.85, 0.85} %{#F8BC65}
\definecolor{funcColor}{rgb}{0.47, 0.6, 0.8}
\definecolor{midgray}{rgb}{0.4, 0.4, 0.3}
\definecolor{indexwhite}{rgb}{0.9, 0.9, 0.9}

\newcommand{\FrameTitle}[2]{\frametitle{#1 \hfill \small\textnormal{\textcolor{midgray}{#2}}}}
\newcommand{\snl}{\vspace*{4pt}\\}
\newcommand{\nl}{\vspace*{1em}\\}
\newcommand{\vsp}{\vspace*{1em}}
\newcommand{\SixPtVspace}{\vspace*{6pt}}

\newcommand{\Quiz}[1]{{\large\textbf{\color{question}{#1}}}}
\newcommand{\QuizQuestion}[1]{{\Large\textbf{\color{question}{#1}}}}

\newcommand{\Matrix}[1]{\textbf{#1}}
\newcommand{\Vector}[1]{\textbf{#1}}
\newcommand{\MMat}[1]{\mathbf{#1}}
\newcommand{\MVec}[1]{\mathbf{#1}}
\newcommand{\Norm}[1]{\lVert {\mathbf{#1}} \rVert}
\newcommand{\NormBeta}[1]{\lVert \beta \mathbf{#1} \rVert}
\newcommand{\NormConst}[1]{\lVert #1 \rVert}
\newcommand{\Trans}{\text{\raisebox{0.5ex}{\scalebox{0.75}{T}}}}
\newcommand{\abs}[1]{\lvert {#1} \rvert}
\newcommand{\Avg}[1]{\mathbf{avg}(\MVec{#1})}
\newcommand{\Std}[1]{\mathbf{std}(\MVec{#1})}
\begin{document}
\title{\center{Applied Linear Algebra\\\vsp Least Squares\\}}

\author{\center{Devi Prasad M \\ Manipal School of Information Sciences \\ Manipal \\}}
\date{\center{September 2026}}
\maketitle


\begin{frame}\FrameTitle{Least Squares}{Unit 1}
For most choices of \Vector{b}, however, there is no $n$-vector
\Vector{x} for which $\MMat{A} \MVec{x} = \MVec{b}$. As a
compromise, we seek an $x$ for which $\MVec{r} = \MMat{A} \MVec{x} - \MVec{b}$,
which we call the \emph{residual} (for the equations $\MMat{A} \MVec{x} - \MVec{b}$),
is as small as possible.

\vsp

We should choose $\MVec{x}$ so as to minimize the norm of the
residual, $\MMat{A} \MVec{x} - \MVec{b}.$

\end{frame}


\begin{frame}\FrameTitle{Least Squares}{Unit 1}
Minimizing the norm of the residual and its square are the same, so we can just
as well minimize
\[
    \lVert \MMat{A} \MVec{x} - \MVec{b} \rVert^2 = \lVert r \rVert^2 = r^2_1 + \cdots + r^2_m,
\]
the sum of squares of the residuals.
\end{frame}
